\input TinLatex.tex
% \usepackage{hyperref}
% \hypersetup{colorlinks,linkcolor={blue},citecolor={blue},urlcolor={red}}  
\setlength{\paperheight}{11in}
\setlength{\paperwidth}{8.5in}
% \bibliographystyle{IEEEtranS} % sorted IEEE style
% \bibliographystyle{IEEEtran} % Đặt kiểu trích dẫn là IEEEtran
\usepackage{multirow}
\usepackage{enumitem}
\usepackage[ruled,linesnumbered]{algorithm2e}
	\SetAlgorithmName{Algorithm}{algorithm}{}
\usepackage{amsmath}

\NewDocumentCommand{\MYref}{O{black}mo}{%
  \begingroup
  \hypersetup{linkcolor=#1}%
  \ref{#2}%
  \IfValueT{#3}{%
    \color{#1}{#3}%
  }%
  \endgroup
}% %NEW COMMAND ALTERNATIVE TO REF

\newcommand\MYreforig[3][black]{\begingroup\hypersetup{linkcolor=#1}\ref{#3}{\color{#1}{#3}}\endgroup}% 

\usepackage{caption} 
\usepackage{cleveref}
\usepackage{color,soul}
\usepackage{pagecolor}
\usepackage{amsthm}
\usepackage{amsfonts}
\usepackage{amssymb}
\usepackage{amsmath}
\newtheorem{assumption}{Assumption}
\newtheorem{remark}{Remark}
\newtheorem{lemma}{Lemma}
\newtheorem{theorem}{Theorem}
\newtheorem{definition}{Definition}
% \newtheorem{proof}{Proof}
\newtheorem{example}{Example}
\usepackage{orcidlink}
\usepackage{algpseudocode}
\usepackage{indentfirst}
\usepackage[T5]{fontenc}
\usepackage{graphicx}
\usepackage[utf8]{inputenc}
\usepackage[vietnamese,english]{babel}
\usepackage{booktabs}
\usepackage{array}






% \renewcommand\hl{}
\begin{document}
\makeatletter	   % `@' is now a normal "letter' for LaTeX
\setcounter{page}{1}
% \renewcommand{\ps@plain}{%
% 	\renewcommand{\@oddhead}{\hfill\begin{tabular}{r}
% 			\small{\emph{Journal of Computer Science and Cybernetics, V.40, N.3 (2025), 1-}} \\
% 			\footnotesize{DOI: 10.15625/1813-9663/19558}
% 		\end{tabular}
% 	}
	
% 	\renewcommand{\@evenhead}{\@oddhead}
% 	\renewcommand{\@oddfoot}{
% 		%\small\hfill{\hskip7.1cm \copyright\ 2024 Vietnam  Academy of Science \& Technology}
% 		\renewcommand{\@evenfoot}{\@oddfoot}
% 	}
% }

%Lenh duoi thay doi do dai cua duong ngang
%\usepackage[flushmargin]{footmisc} - PHẢI MỞ LỆNH NÀY
\def\footnoterule{\kern-3\p@
	\hrule \@width 1.32in \kern 2.6\p@} % the \hrule is .4pt high
%Thay doi indent cua footnote
%\renewcommand{\footnotesep}{2cm}
% \renewcommand{\footnotemargin}{0em}
\newcommand{\fn}[1]{\footnotetext{\hspace{-6mm}#1}}
\setlength{\skip\footins}{8mm}

\makeatother   % `@' is restored as a "non-letter" character 
    
%%%%%%%%%%%%%%%%%%%%%%%%
%Title & Authors
%%%%%%%%%%%%%%%%%%%%%%%%

%=========================================================================
\title{ENHANCING MULTILINGUAL VISION-LANGUAGE ALIGMENT WITH TEXT DATA AUGMENTATION: A CASE STUDY OF VIETNAMESE mCLIP}

\author{
	{\cn PHUONG THAO HOANG, THI HOAI PHAN, THI-LAN LE, THI NGOC DIEP DO*}
	\vskip.5cm
	{
    \it SigM Laboratory, School of Electrical and Electronic Engineering \\ Hanoi University of Science and Technology\\
        \it No. 1 Dai Co Viet, Bach Mai, Ha Noi City, Viet Nam\\
\fn{*Corresponding author.}
\fn{\hspace{1.7mm}{\it E-mail addresses}:
\href{mailto:Thao.HP224275@sis.hust.edu.vn}{Thao.HP224275@sis.hust.edu.vn} (P.T HOANG); 
\href{mailto:anonymous@ioit.ac.vn}{hoai.pt220031d@sis.hust.edu.vn} (THI HOAI PHAN);
\href{mailto:anonymous@ioit.ac.vn}{lan.lethi1@hust.edu.vn} (THI LAN LE);
\href{mailto:anonymous@ioit.ac.vn}{diep.dothingoc@hust.edu.vn} (THI NGOC DIEP DO).
}}}
\maketitle
\renewcommand\refname{\normalsize \centerline{ REFERENCES}}
% Set to use the "plain" pagestyle
\pagestyle{plain}
\pagestyle{myheadings}
\markboth{\footnotesize \chu \uppercase{HOANG PHUONG THAO} {\it et al.}}
{\footnotesize  \chu \uppercase{ENHANCING MULTILINGUAL VISION-LANGUAGE ALIGMENT WITH TEXT AUGMENTATION }}% Maximun 8 words 
{	\vspace{-.4cm}
	%\hspace{13.25cm}{\includegraphics[scale=.2]{crosscheck.pdf}}%\vspace{-.5cm}
	\hspace{12.50cm}{\includegraphics[scale=1.2]{crossref.png}
	}
}
%=========================================================================

{\cn
%=========================================================================
\begin{abstract}{ Large-scale vision-language pretrained (VLP) models such as CLIP have demonstrated strong performance in cross-modal tasks, but their effectiveness remains limited for low-resource languages due to the dominance of English training data. This work investigates a data-efficient strategy for adapting multilingual CLIP (mCLIP) to Vietnamese without modifying pretrained visual encoders. Rather than proposing a new architecture, we focus on strengthening target-language supervision through translated captions and text data augmentation. We first construct a Vietnamese version of the Multi30k dataset by translating English captions and enriching them with augmentation techniques including random insertion, random deletion, random swap, and back-translation to increase linguistic diversity. The adapted model is evaluated on both the translated Vietnamese Multi30k dataset and UIT-ViIC, a native Vietnamese image-caption dataset. Experiments are conducted on cross-modal retrieval tasks to evaluate the quality of vision-language alignment.} Results show that text augmentation consistently improves retrieval performance across both datasets. Mean recall (MeanR) increases from 72.3 to 73.3 on Vietnamese Multi30k and from 45.9 to 47.2 on UIT-ViIC. Notably, image-to-text Recall@1 improves from 49.3 to 51.0 and from 20.2 to 22.3, respectively. To assess robustness beyond the specific-domain bias of UIT-ViIC, we further evaluate the adapted model on another Vietnamese dataset containing more general images, providing an assessment of cross-domain generalization. Although performance is slightly lower than that on the in-domain UIT-ViIC dataset, the model retains reasonable Recall@K scores.
}These findings demonstrate that effective multilingual transfer to Vietnamese can be achieved through modest text-side augmentation and parameter-efficient fine-tuning. 
\vn
\end{abstract}
\textbf{Keyword.} Contrastive learning, cross-lingual transfer, data augmentation, image–text retrieval, mclip, multilingual learning, vietnamese language processing, vision–language models.
}


%%%%%%%%%%%%%%%%%%%%%%%%
%Main text
%%%%%%%%%%%%%%%%%%%%%%%%

%=========================================================================
\section{INTRODUCTION}

Large-scale vision--language pretrained (VLP) models such as CLIP~\cite{radford2021learningtransferablevisualmodels} have demonstrated strong generalization across a wide range of multimodal tasks. Trained contrastively on hundreds of millions of image--text pairs, CLIP learns a shared embedding space that tightly aligns visual and textual representations, while its dual-encoder architecture enables efficient large-scale retrieval. Despite its strong performance, CLIP remains largely English-centric due to the dominance of English data during pretraining. Consequently, its effectiveness degrades when applied to low-resource languages such as Vietnamese, where large-scale aligned image--text corpora are scarce and linguistic divergence from English is substantial. Existing approaches for extending CLIP to multilingual settings typically involve fine-tuning the English text encoder using parallel corpora or replacing it with multilingual transformers such as mBERT~\cite{pires2019multilingualmultilingualbert} or XLM-R~\cite{conneau2020unsupervisedcrosslingualrepresentationlearning}. Although these strategies introduce multilingual capability, they are often computationally expensive and risk disrupting CLIP’s well-optimized embedding geometry, potentially reducing retrieval efficiency and scalability. These limitations motivate parameter-efficient transfer strategies that preserve pretrained components while extending language coverage. Recent parameter-efficient fine-tuning techniques such as Adapter-tuning and Low-Rank Adaptation (LoRA) \cite{hu2021loralowrankadaptationlarge} have also been widely explored for adapting large pretrained models. These approaches introduce lightweight trainable modules while keeping most backbone parameters frozen, enabling efficient adaptation without retraining the entire model.

The mCLIP framework~\cite{chen-etal-2023-mclip} addresses this challenge by decoupling multilingual language modeling from multimodal alignment. Instead of retraining CLIP, mCLIP freezes the original CLIP image and English text encoders and aligns a multilingual text encoder to the CLIP embedding space via lightweight projection layers optimized with Triangle Cross-modal Knowledge Distillation. By updating only a small fraction of parameters, mCLIP enables efficient adaptation to low-resource languages while preserving pretrained visual semantics. Importantly, visual representations learned by CLIP are grounded in image content rather than tied to a specific language. Core visual concepts such as objects, actions, attributes, and spatial relations remain consistent across linguistic descriptions. While captions may differ between English and Vietnamese, the underlying visual semantics encoded by the image backbone are largely language-independent. Consequently, cross-lingual discrepancies primarily arise in textual representation rather than visual features. This observation suggests that, in low-resource scenarios, effective adaptation may benefit more from strengthening target-language textual supervision than from retraining the visual encoder. Vietnamese provides a compelling test case due to the scarcity of large-scale Vietnamese image--text datasets. To investigate adaptation under both synthetic and native supervision, we consider two complementary datasets. First, we construct a Vietnamese version of Multi30k~\cite{elliott2016multi30k} by translating English captions and enriching them with controlled text augmentation techniques to increase lexical and syntactic diversity. Second, we evaluate on UIT-ViIC~\cite{lam2020uitviicdatasetevaluationvietnamese}, a native Vietnamese image--caption dataset with human-written annotations. Although UIT-ViIC is restricted to the sport-ball domain, it provides a realistic evaluation setting for adaptation under authentic Vietnamese linguistic conditions. To further assess cross-domain generalization, we additionally evaluate the adapted model on a Vietnamese dataset containing non-sports imagery. This evaluation provides an out-of-distribution setting for examining whether the learned cross-modal alignment transfers beyond the sports-specific context of UIT-ViIC.


In this work, we adapt mCLIP to Vietnamese and systematically investigate the role of text data augmentation in improving cross-modal alignment under both translated-caption and native-caption scenarios. Rather than proposing a new architecture, we adopt a data-centric perspective and establish a practical and reproducible adaptation pipeline that combines modest data expansion with parameter-efficient fine-tuning. Our objective is to analyze whether controlled linguistic diversification alone can enhance embedding quality without increasing model complexity.

The main contributions of this work are summarized as follows:

\begin{itemize}
\item We investigate the adaptation of the multilingual vision--language model mCLIP to Vietnamese under low-resource conditions while preserving the pretrained visual encoder.

\item We construct a Vietnamese version of the Multi30k dataset by translating English captions into Vietnamese, thereby creating a translated benchmark for Vietnamese vision--language adaptation.

\item We apply text data augmentation to both the translated Vietnamese Multi30k training set and the native Vietnamese UIT-ViIC training set in order to increase linguistic diversity and strengthen target-language supervision.

\item We conduct experiments on both Vietnamese Multi30k and UIT-ViIC for cross-modal retrieval, and provide an empirical analysis showing that text-side augmentation consistently improves cross-modal alignment for Vietnamese.
\end{itemize}


Experimental results demonstrate that text augmentation consistently improves Vietnamese image--text retrieval performance on both datasets, with particularly notable gains at stricter metrics such as Recall@1, indicating enhanced fine-grained cross-modal alignment.

The remainder of this paper is organized as follows. Section~\ref{sec:related} reviews related work on multilingual vision--language pretraining. Section~\ref{sec:method} describes the proposed adaptation and optimization pipeline. Section~\ref{sec:dataset} presents dataset construction and augmentation strategies. Section~\ref{sec:experiments} details the experimental settings and evaluation protocols. Section~\ref{sec:results} reports the results and analysis, and Section~\ref{sec:conclusion} concludes the paper.



\section{RELATED WORK}
\label{sec:related}

Vision--language pretraining (VLP) learns transferable multimodal representations from large-scale image--text pairs. Contrastive dual-encoder models such as CLIP~\cite{radford2021learningtransferablevisualmodels} and ALIGN~\cite{jia2021scalingvisualvisionlanguagerepresentation} align images and captions in a shared embedding space, enabling efficient retrieval by independently encoding each modality and applying similarity search. Despite strong zero-shot generalization, these models inherit a pronounced English-centric bias because their web-scale corpora are overwhelmingly English, leading to severe degradation when applied to low-resource languages such as Vietnamese.

\subsection{Multilingual VLP via Synthetic Supervision}
A common strategy for multilingual VLP is to introduce synthetic multilingual supervision, typically by translating English captions or injecting multilingual textual objectives. M3P~\cite{ni2021m3plearninguniversalrepresentations} combines multilingual masked language modeling with cross-modal pretraining objectives to learn more language-agnostic representations, improving cross-lingual transfer compared to monolingual baselines but still requiring large multilingual corpora and heavy pretraining. UC2~\cite{zhou2021uc2universalcrosslingualcrossmodal} further explores universal cross-lingual cross-modal alignment using translated captions and multilingual objectives; however, it often employs single-stream or tightly coupled architectures where image and text tokens interact within a unified Transformer, increasing inference cost and reducing scalability for retrieval-oriented deployments.

MURAL~\cite{jain2021muralmultimodalmultitaskretrieval} scales multilingual retrieval by combining translated image--text pairs with multilingual text-only contrastive learning. Its key insight is to strengthen multilingual sentence representations via large parallel corpora and then anchor them to visual concepts through translated supervision. While effective at broad language coverage, such scaling-based methods typically demand substantial compute and extensive multilingual corpora, which are difficult to obtain and reproduce for truly low-resource languages.

\subsection{Multilingual CLIP Adaptation}
Beyond training multilingual VLP models from scratch, another line of work adapts pretrained CLIP models to multilingual settings. A straightforward approach replaces CLIP's English text encoder with multilingual transformers (e.g., mBERT~\cite{pires2019multilingualmultilingualbert} or XLM-R) and fine-tunes the resulting system on multilingual data. Although this enables multilingual inputs, full encoder replacement is computationally expensive and risks distorting CLIP’s well-structured embedding geometry, potentially harming English performance and weakening the efficiency advantages of dual-encoder retrieval.

In contrast, mCLIP~\cite{Chen_2023} proposes a parameter-efficient cross-lingual transfer framework that preserves CLIP’s pretrained image and English text encoders and aligns multilingual text representations via lightweight projection layers. Triangle Cross-modal Knowledge Distillation (TriKD) enforces consistency among image--English text, image--multilingual text, and English text--multilingual text similarities, enabling transfer to new languages while maintaining retrieval efficiency. Our work follows this parameter-efficient philosophy but emphasizes a complementary direction: we study how translation and, crucially, text augmentation can strengthen Vietnamese supervision and improve robustness under limited data.

\begin{figure}[t]
    \centering
    \includegraphics[width=0.95\columnwidth]{block_diagram.jpg}
    \caption{Overview of the proposed adaptation pipeline.}
    %\textcolor{red}{vẽ lại phần mCLIP to hơn, các khối phía trên bé đi }
    \label{fig:pipeline}
\end{figure}



\subsection{Advanced Vision--Language Models}
Recent VLP models go beyond CLIP-style contrastive training by incorporating stronger objectives or generative components. BLIP~\cite{li2022blipbootstrappinglanguageimagepretraining} couples contrastive alignment with image--text matching and captioning losses, and introduces bootstrapping mechanisms to improve pretraining data quality; this generally improves cross-modal understanding but relies on additional training stages and large-scale data curation. PaLI~\cite{chen2023palijointlyscaledmultilinguallanguageimage} demonstrates that jointly scaling a large multilingual language model with vision inputs and multi-task supervision yields strong multilingual transfer, yet it requires massive training budgets and large multilingual corpora, making it less suitable for lightweight adaptation scenarios. SigLIP~\cite{tschannen2025siglip} revisits contrastive learning with improved objectives and training recipes that can enhance retrieval performance and training stability; however, it is typically studied in large-scale settings and does not directly address low-resource language adaptation when paired target-language image--text data are scarce. 

\subsection{Vietnamese Multimodal Resources}
Vietnamese multimodal learning remains under-explored compared to high-resource languages, primarily due to limited availability of large-scale Vietnamese image--text datasets. There are few of publicly reported datasets with human-written Vietnamese captions. The scale and domain coverage are relatively constrained also. This motivates data-centric adaptation strategies that maximize the utility of existing resources. Rather than proposing a new architecture, we investigate whether a modest translated corpus enriched with text augmentation can effectively adapt mCLIP\cite{chen-etal-2023-mclip} to Vietnamese, and we further validate that augmentation remains beneficial on native Vietnamese captions datasets.


\section{METHODOLOGY}
\label{sec:method}
%\subsection{mCLIP Model Overview}

mCLIP~\cite{chen-etal-2023-mclip} follows the dual-encoder architecture of CLIP, consisting of a pretrained CLIP ViT image encoder and a multilingual text encoder (MTE) initialized from XLM-R, followed by a projection module. The pretrained CLIP image encoder remains aligned with the English text encoder, while the multilingual encoder is trained to map target-language text into the same embedding space. Triangle Cross-modal Knowledge Distillation (TriKD) preserves the learned image–English alignment while transferring multilingual capability.

% Figure~\ref{fig:pipeline} illustrates our adaptation for a new target language. The CLIP image encoder is kept frozen to retain strong visual representations learned from large-scale English pretraining. Only the projection layers, including the shared CLIP projector and the X-Projector, are updated. This parameter-efficient design aligns target-language textual embeddings with the fixed visual space without disrupting pretrained representations.

Figure~\ref{fig:pipeline} illustrates the overall pipeline of the proposed Vietnamese adaptation framework. Starting from the original English captions of the Multi30k dataset, captions are first translated into Vietnamese using a neural machine translation system. To increase linguistic diversity and mitigate structural artifacts introduced by translation, several text augmentation techniques are applied, including random insertion, random deletion, random swap, and back-translation. The augmented captions are then combined with the translated captions and processed through standard preprocessing steps to produce a training-ready corpus. During adaptation, images are encoded using the frozen CLIP image encoder, while Vietnamese captions are encoded by the multilingual text encoder. The resulting text representations are projected into the shared CLIP embedding space through lightweight projection modules. To preserve the pretrained multimodal representation, the CLIP image encoder and the multilingual text encoder remain frozen during training. Only the projection layers (including the CLIP projector and the X-projector) are updated, alternative parameter-efficient approaches such as LoRA could also be applied by injecting low-rank trainable matrices into attention layers of the multilingual encoder. However, in this work we primarily focus on projection-layer adaptation because it directly aligns textual representations with the fixed CLIP embedding space while preserving the internal structure of the pretrained encoder.



\subsection{Image--Text Alignment for a Target Language}

Given an image–text dataset in the target language, images are encoded by the frozen CLIP image encoder, while captions are processed by the multilingual text encoder and projected into CLIP’s embedding space. Alignment is learned via a contrastive objective that pulls matched image–text pairs closer and pushes mismatched pairs apart within a minibatch.

Let $\{(I_i, X_i)\}_{i=1}^{N}$ denote a minibatch of image–text pairs. The normalized image embedding is $h_I^i \in \mathbb{R}^d$, and the normalized caption embedding is $h_X^i \in \mathbb{R}^d$. The directional contrastive loss is

\[
\ell(x, y)
=
-
\frac{1}{N}
\sum_{i=1}^{N}
\log
\frac{\exp(x_i^\top y_i / \tau)}
{\sum_{j=1}^{N} \exp(x_i^\top y_j / \tau)},
\]

where $\tau$ is a temperature parameter.

\paragraph{Image--Text Contrastive Loss}
The symmetric image–text contrastive loss is defined as
\[
\mathcal{L}_{\mathrm{ITC}}
=
\frac{1}{2}
\big[
\ell(h_I, h_X)
+
\ell(h_X, h_I)
\big].
\]

\paragraph{TriKD and Target-Only Training}
When paired English captions are available, mCLIP additionally applies a text–text contrastive loss between the frozen CLIP English text embedding $h_T^i$ and the target-language embedding $h_X^i$:
\[
\mathcal{L}_{\mathrm{TTC}}
=
\frac{1}{2}
\big[
\ell(h_T, h_X)
+
\ell(h_X, h_T)
\big],
\]
leading to the full TriKD objective
\[
\mathcal{L}_{\mathrm{TriKD}}
=
\mathcal{L}_{\mathrm{ITC}}
+
\lambda \mathcal{L}_{\mathrm{TTC}}.
\]

For non-multilingual datasets which contains only target-language captions, the text–text term cannot be computed. In this case, training reduces to
\[
\mathcal{L}_{\mathrm{train}} = \mathcal{L}_{\mathrm{ITC}}.
\]

This setting can be viewed as a simplified form of TriKD, focusing solely on image–target alignment while preserving the pretrained visual space. Since only lightweight projection layers (approximately 3\% of model parameters) are updated, adaptation remains efficient and stable.

\subsection{Target-language dataset construction and text augmentation}
To construct a target-language version of an image–English dataset, we translate all English captions using a high-quality neural machine translation system. Each caption is translated independently to preserve the original one-to-many relationship between images and their descriptions, resulting in a parallel corpus with translated captions in the target language. Although translation retains semantic content, it may introduce artifacts that affect downstream alignment.

Machine-translated captions often inherit rigid syntactic structures from English, such as fixed subject–verb–object order, literal idiomatic expressions, and limited lexical diversity. These artifacts reduce linguistic variation and may bias the multilingual text encoder toward unnatural surface forms. Throughout preprocessing, image identifiers and official dataset splits are preserved to ensure compatibility with the original structure while providing a complete target-language caption layer for adaptation.

To enhance linguistic diversity and mitigate structural rigidity, we apply controlled text augmentation techniques that introduce lexical and syntactic variation while preserving semantic meaning~\cite{shorten2021text, wei2019eda}. These augmentations expand the range of surface realizations associated with each image and encourage the multilingual text encoder to learn representations that are robust to natural language variation. An example is illustrated in Table~\ref{tab:augmentation}.

We employ four augmentation strategies.

\textit{Random deletion} removes tokens independently with probability $p$, ensuring that at least one token remains. This generates shorter paraphrases that emphasize essential content words while discarding less informative modifiers. Exposure to such variants encourages the model to rely on core semantic cues rather than superficial syntactic structure.

\textit{Random insertion} introduces additional tokens at plausible positions within the caption. Inserted tokens are selected from a set of common descriptive words or short phrases that refine nearby content. This strategy increases lexical diversity and introduces mild stylistic variation that better reflects naturally occurring descriptions.

\textit{Random swap} exchanges the positions of two tokens or short phrases, generating syntactic variants without altering meaning. Since many languages permit flexible word order, training on swapped variants encourages invariance to local reordering and reduces sensitivity to strict syntactic patterns introduced by machine translation.

\textit{Back-translation}~\cite{corbeil2020bet} produces deeper paraphrastic variants by translating each caption into an intermediate language and then back into the target language. Compared to token-level edits, back-translation introduces broader lexical and structural variation while preserving the underlying semantics. This global augmentation complements local perturbations and contributes more diverse training signals.

By combining translated captions with augmentation techniques, we construct a linguistically enriched image–text corpus that increases diversity in the target language. The augmented data enables the multilingual text encoder to learn more robust representations that align effectively with CLIP’s fixed visual embedding space without modifying the pretrained image backbone.

All captions are processed using the \texttt{fairseq} pipeline. Frequency-based vocabularies are built, low-frequency tokens are mapped to \texttt{<unk>}, and the text is binarized for efficient training. Augmented captions are treated as independent samples and integrated with the original translations, producing a training-ready dataset for parameter-efficient mCLIP adaptation.
\begin{table}[ht]
\centering
\caption{Examples of text augmentation strategies applied to a Vietnamese caption and English caption}
\label{tab:augmentation}
\begin{tabular}{|l|p{5cm}|p{5cm}|}
\hline
\textbf{Augmentation} & \textbf{Vietnamese example} & \textbf{English example} \\
\hline
Original &
Hai vận động viên tennis đang luyện tập trên sân &
Two tennis players are practicing on the court \\
\hline
Random deletion &
vận động viên tennis đang luyện tập ở trên sân &
Two players practicing on the \\
\hline
Random insertion &
Hai vận động viên tennis đang luyện luyện tập ở trên sân &
Two tennis players are practicing on the participant court \\
\hline
Random swap &
Hai vận viên tennis đang luyện tập ở động sân &
Two the players are practicing on tennis court \\
\hline
Back-translation &
Hai tay vợt đang tập luyện trên sân &
Two tennis players are training on the court \\
\hline
\end{tabular}
\end{table}
\section{DATASETS}
\label{sec:dataset}

\subsection{Vietnamese Multi30k}

Multi30k~\cite{elliott2016multi30k} is a widely used benchmark for multilingual and cross-modal learning, extending Flickr30k dataset \cite{plummer2016flickr30kentitiescollectingregiontophrase} dataset with professionally written English captions. It contains approximately 31{,}000 images depicting everyday scenes, each paired with five human-authored English descriptions. We adopt the official train, validation, and test splits without modification to ensure comparability with prior work.

To construct a Vietnamese counterpart, all English captions from the original Multi30k dataset are translated into Vietnamese using a transformer-based neural machine translation (NMT) API\cite{google_translate_api}. Each caption is translated independently to preserve the original one-to-many image–caption correspondence. The resulting Vietnamese training split contains approximately 29K captions, while the validation and test splits contain approximately 1K captions each.

Although translation preserves semantic content, machine-translated captions often exhibit rigid syntactic patterns and limited lexical variation, frequently reflecting English word order. Such artifacts may reduce linguistic diversity and limit the robustness of learned representations.

To address this issue, we apply text augmentation exclusively to the Vietnamese training split using four techniques: random deletion, random insertion, random swap, and back-translation. For each translated caption, two augmented variants are generated, expanding the training corpus to approximately 145K captions. The validation and test splits remain unchanged to ensure fair evaluation.

This translated dataset allows us to investigate Vietnamese adaptation under synthetic supervision, where linguistic diversity is constrained by translation artifacts.


\subsection{UIT-ViIC}

UIT-ViIC~\cite{lam2020uitviicdatasetevaluationvietnamese} is a manually annotated Vietnamese image--caption dataset constructed from a subset of MS-COCO images. It contains 3{,}850 images from the sport-ball domain, each paired with five human-written Vietnamese captions, resulting in 19{,}250 captions in total.

Unlike translated corpora, UIT-ViIC captions are written directly in Vietnamese and therefore reflect natural lexical choices, flexible word order, and idiomatic expressions. This makes the dataset a realistic benchmark for evaluating cross-modal alignment under native-language supervision.

It is important to note that UIT-ViIC is restricted to a single semantic domain. Vocabulary related to sports activities and equipment appears frequently, resulting in a relatively concentrated lexical distribution and reduced cross-domain variability. While such domain specificity enables controlled evaluation within a consistent visual context, it may limit the generalizability of conclusions to broader everyday imagery. Consequently, results on UIT-ViIC should be interpreted as evidence of adaptation under domain-constrained native supervision rather than as a fully domain-diverse assessment.

In addition, the relatively small scale of the dataset presents a challenging low-resource adaptation scenario. The original training split contains approximately 13.4K captions. To alleviate data sparsity, we apply the same four augmentation techniques used for Vietnamese Multi30k to the training set only, expanding it to approximately 67K captions. Validation and test sets remain unchanged to ensure fair evaluation.

By jointly evaluating Vietnamese Multi30k and UIT-ViIC, we examine adaptation under both translated supervision and domain-constrained native supervision.



Table \ref{tab:dataset-stats} summarizes the datasets statistic before and after augmentation.
\begin{table}[ht]
\centering
\caption{Statistics of Vietnamese data before and after augmentation}
\label{tab:dataset-stats}
\begin{tabular}{|c|c|c|c|c|c|}
\hline
\textbf{Dataset} & \textbf{Train data}  & \textbf{Augmented} & \textbf{Valid data} & \textbf{Test data} \\
\hline
Vietnamese Multi30k & 29K & +116K  & 1014 & 1000 \\ 
\hline
UIT-ViIC & 13.4K & +53.6K & 4620 & 1155 \\ 
\hline
\end{tabular}
\end{table}
\subsection{KTVIC (Out-of-Domain Evaluation)}


To evaluate cross-domain generalization, we additionally test the adapted model on the KTVIC dataset \cite{pham2024ktvicvietnameseimagecaptioning}, which contains Vietnamese captions describing general everyday scenes rather than sports-related imagery. The dataset consists of 3,769 images in the training split and 558 images in the test split. Since our objective is to evaluate cross-domain robustness rather than perform additional training, we use only the test split for retrieval evaluation. To construct a one-to-one image–caption retrieval setting, we select the first caption associated with each image, resulting in 558 image–caption pairs for evaluation. Unlike UIT-ViIC, whose images predominantly depict sport-ball activities, KTVIC provides a more diverse set of visual contexts and vocabulary. This dataset therefore serves as an out-of-distribution evaluation for models trained or fine-tuned on the sports-domain UIT-ViIC dataset. By measuring retrieval performance on KTVIC without additional fine-tuning, we assess whether the learned cross-modal alignment remains effective when applied to more general images in Vietnamese.

\section{EXPERIMENTS}
\label{sec:experiments}
\subsection{Experimental Setup and Ablation Configuration}

All models are initialized from a publicly released mCLIP checkpoint, comprising a pretrained CLIP ViT image encoder and a multilingual text encoder aligned to the CLIP embedding space. During adaptation, only the lightweight projection layers are updated, while all backbone encoders remain frozen. This parameter-efficient design preserves the pretrained visual representation while enabling effective cross-lingual alignment.

Experiments are implemented using the \textit{fairseq} framework and conducted on a single NVIDIA GPU. We use the official train, validation, and test splits of Vietnamese Multi30k and UIT-ViIC.

To evaluate the contribution of each augmentation technique, we perform controlled ablation experiments on Vietnamese Multi30k. Random deletion, random insertion, random swap, and back-translation are applied individually under identical training settings. The baseline corresponds to training without augmentation. Batch size, learning rate, training steps, and early stopping criteria are kept constant across all variants to ensure fair comparison and isolate the impact of each augmentation method on retrieval performance.


\subsection{Training Details}

We fine-tune mCLIP using the symmetric image--text contrastive loss described in Section~III. Mini-batches consist of image--caption pairs sampled from the target-language training data. When augmentation is enabled, augmented captions are treated as independent samples.

Optimization is performed using Adam \cite{jais2019adam} with a learning rate selected from $10^{-5}$ to $10^{-4}$ based on validation performance. The temperature parameter $\tau$ in the contrastive loss follows the default configuration of mCLIP. Early stopping is applied based on validation MeanR to prevent overfitting.

For Vietnamese Multi30k, text augmentation increases the training set from approximately 29K translated captions to 145K captions. For UIT-ViIC, augmentation expands the training data from 13.4K to approximately 67K captions. All other hyperparameters are kept identical across experiments to ensure fair comparison. Table~\ref{tab:training_params} summarizes the main experimental hyperparameters used for adapting the mCLIP model to Vietnamese.

\begin{table}[ht]
\centering
\caption{Experimental parameter settings for mCLIP adaptation.}
\label{tab:training_params}
\scriptsize
\renewcommand{\arraystretch}{1.2}
\begin{tabular}{|l|l|}
\hline
\textbf{Parameter} & \textbf{Value} \\
\hline
Image encoder & CLIP ViT (frozen) \\
\hline
Text encoder & XLM-R multilingual encoder (frozen) \\
\hline
Trainable modules & Projection layers (CLIP projector, X-Projector) \\
\hline
Optimizer & Adam \\
\hline
Learning rate & $1 \times 10^{-5}$ -- $1 \times 10^{-4}$ \\
\hline
Batch size & 32 \\
\hline
Temperature ($\tau$) & Default mCLIP configuration \\
\hline
Training framework & fairseq \\
\hline
Hardware & Single NVIDIA GPU \\
\hline
Embedding normalization & $\ell_2$ normalization \\
\hline
Retrieval backend & FAISS IndexFlatL2 \\
\hline
Evaluation metrics & Recall@1, Recall@5, Recall@10, MeanR \\
\hline
\end{tabular}
\end{table}
\subsection{Evaluation Protocol}

We evaluate cross-modal retrieval performance in both image-to-text (i2t) and text-to-image (t2i) directions. Performance is measured using Recall@K with $K \in \{1,5,10\}$. A query is considered correct if at least one ground-truth match appears within the top-$K$ retrieved results. We additionally report mean recall (MeanR), computed as the average of the six Recall@K scores, as an overall evaluation metric. At inference time, all image and text embeddings are $\ell_2$-normalized. Nearest-neighbor search is performed using FAISS with an IndexFlatL2 index. For datasets with multiple captions per image, retrieved caption indices are mapped back to image identifiers and duplicates are removed while preserving ranking order. We focus on bidirectional retrieval as a foundational evaluation of cross-modal embedding quality. Since mCLIP adopts a dual-encoder architecture optimized for similarity-based alignment, cross-modal retrieval provides a direct and principled evaluation of the shared embedding space. Retrieval performance therefore reflects the effectiveness of multilingual alignment under the parameter-efficient adaptation setting considered in this study.

In addition to in-domain evaluation on Vietnamese Multi30k and UIT-ViIC, we conduct an out-of-distribution test using the KTVIC dataset \cite{pham2024ktvicvietnameseimagecaptioning}. The model is evaluated without additional fine-tuning in order to measure the robustness of the learned cross-modal alignment when applied to images from a different semantic domain.



\section{RESULTS AND ANALYSIS}
\label{sec:results}
This section reports cross-modal retrieval results on Vietnamese Multi30k and UIT-ViIC, with a detailed analysis of the contribution of text data augmentation. Since mCLIP follows a dual-encoder architecture optimized for similarity-based alignment, retrieval performance directly reflects the quality and integrity of the shared embedding space. Evaluating retrieval therefore provides a principled assessment of multilingual alignment before extending to higher-level multimodal reasoning tasks.

\subsection{Results on Vietnamese Multi30k}

Table~\ref{tab:multi30k-results} reports cross-modal retrieval performance on the Vietnamese Multi30k dataset. Fine-tuning mCLIP on machine-translated Vietnamese captions already yields strong results, indicating that translated supervision is sufficient to align Vietnamese textual representations with the frozen CLIP visual encoder. Without augmentation, the model achieves a MeanR of 72.3.

As shown in Table~\ref{tab:multi30k-results}, incorporating text augmentation consistently improves performance in both image-to-text (i2t) and text-to-image (t2i) retrieval. The largest gains are observed at Recall@1 in both directions, increasing from 49.3 to 51.0 for i2t and from 49.9 to 51.2 for t2i. Improvements are also observed at higher recall levels (R@5 and R@10). Overall, MeanR increases from 72.3 to 73.3 after applying the full augmentation pipeline. These results indicate that modest textual diversification can strengthen cross-modal alignment without modifying the visual backbone.

\begin{table}[ht]
\centering
\caption{Cross-modal retrieval results on Vietnamese Multi30k}
\label{tab:multi30k-results}
\begin{tabular}{|l|c|c|c|}
\hline
\textbf{Model} & \textbf{R@1} & \textbf{R@5} & \textbf{R@10} \\
\hline
\multicolumn{4}{|c|}{\textit{Image-to-Text (i2t)}} \\
\hline
mCLIP & 49.3 & 80.1 & 86.7 \\
\hline
mCLIP (+ Aug.) & \textbf{51.0} & \textbf{80.5} & \textbf{87.6} \\
\hline
\multicolumn{4}{|c|}{\textit{Text-to-Image (t2i)}} \\
\hline
mCLIP & 49.9 & 80.4 & 87.7 \\
\hline
mCLIP (+ Aug.) & \textbf{51.2} & \textbf{80.7} & \textbf{88.9} \\
\hline
\end{tabular}
\end{table}


\subsection{Results on UIT-ViIC}

Table~\ref{tab:uit} presents cross-modal retrieval results on the UIT-ViIC dataset. Overall retrieval performance is lower than on Vietnamese Multi30k due to the greater linguistic diversity of native Vietnamese captions and the smaller dataset size. The domain-specific nature of UIT-ViIC, which focuses on sport-related imagery, may also influence retrieval dynamics because of concentrated vocabulary and repetitive visual patterns. Without augmentation, the model achieves a MeanR of 45.9.

As shown in Table~\ref{tab:uit}, text augmentation consistently improves retrieval accuracy across most metrics. In the image-to-text (i2t) direction, Recall@1 increases from 20.2 to 22.3, with additional improvements at R@5 and R@10. In the text-to-image (t2i) direction, gains are observed at R@5 and R@10, while Recall@1 remains comparable. Although the absolute improvement in MeanR is modest (from 45.9 to 47.2), the consistent gains—particularly at stricter metrics such as Recall@1—indicate enhanced fine-grained cross-modal alignment.

These findings suggest that textual diversification improves robustness in native-caption settings, even under domain-constrained conditions.

\begin{table}[ht]
\centering
\caption{Cross-modal retrieval results on the UIT-ViIC dataset}
\label{tab:uit}
\begin{tabular}{|l|c|c|c|}
\hline
\textbf{Model} & \textbf{R@1} & \textbf{R@5} & \textbf{R@10} \\
\hline
\multicolumn{4}{|c|}{\textit{Image-to-Text (i2t)}} \\
\hline
mCLIP & 20.2 & 44.8 & 59.9 \\
\hline
mCLIP (+ Aug.) & \textbf{22.3} & \textbf{46.9} & \textbf{63.1} \\
\hline
\multicolumn{4}{|c|}{\textit{Text-to-Image (t2i)}} \\
\hline
mCLIP & 23.8 & 54.1 & 72.3 \\
\hline
mCLIP (+ Aug.) & 23.4 & \textbf{55.0} & \textbf{72.7} \\
\hline
\end{tabular}
\end{table}

\subsection{Cross-domain evaluation on KTVIC}


To examine whether the learned cross-modal alignment generalizes beyond the sports-domain bias of UIT-ViIC, we evaluate the adapted model on the KTVIC dataset, which contains non-sports images and captions describing general everyday scenes. The KTVIC dataset consists of 558 images in the test split, which is substantially smaller than the UIT-ViIC test set. In our evaluation, we use only the KTVIC test split and construct a one-to-one retrieval setting by selecting the first caption for each image. Table~\ref{tab:ktvic} reports retrieval performance on KTVIC. Although performance is slightly lower than on the in-domain UIT-ViIC dataset, the model retains reasonable Recall@K scores. This result suggests that the learned embedding space preserves useful semantic alignment even under cross-domain conditions. The findings indicate that the parameter-efficient adaptation combined with text augmentation improves model robustness and prevents overfitting to the sports-specific vocabulary and visual patterns present in UIT-ViIC. In the next part, we focus on analyzing the experiment results on UIT-ViIC dataset. 

\begin{table}[ht]
\centering
\caption{Out-of-domain retrieval results on KTVIC using the UIT-ViIC-trained model.}
\label{tab:ktvic}
\begin{tabular}{|l|c|c|c|}
\hline
\textbf{Model} & \textbf{R@1} & \textbf{R@5} & \textbf{R@10} \\
\hline
\multicolumn{4}{|c|}{\textit{Image-to-Text (i2t)}} \\
\hline
mCLIP & 23.7 & 47.0 & 59.1 \\
\hline
\multicolumn{4}{|c|}{\textit{Text-to-Image (t2i)}} \\
\hline
mCLIP & 20.6 & 44.8 & 56.3 \\
\hline
\end{tabular}
\end{table}

\subsection{Ablation Study on Text Augmentation Techniques}

To evaluate the contribution of each augmentation technique, we conduct an ablation study on Vietnamese Multi30k by applying random deletion, random insertion, random swap, and back-translation individually while keeping all other settings fixed. The detailed results are presented in Table \ref{tab:ablation}, reporting R@1, R@5, and R@10 for both image-to-text (i2t) and text-to-image (t2i) retrieval.

As shown in Table \ref{tab:ablation}, all augmentation methods improve performance over the baseline across both retrieval directions and all Recall@K metrics. Among the four individual strategies, random deletion achieves the strongest and most consistent gains on R@1, R@5, and R@10 in both i2t and t2i. This indicates that removing non-essential tokens helps the model focus on core semantic content and enhances cross-modal alignment.

While the other techniques also provide improvements, their gains are comparatively smaller. Combining all four strategies yields the best overall performance, but the ablation results highlight random deletion as the most effective single augmentation method.

This systematic comparison clarifies the empirical contribution of this study, demonstrating that structured data diversification alone can measurably strengthen cross-modal alignment without introducing architectural modifications.



\begin{table*}[ht]
\centering
\caption{Ablation study of individual text augmentation techniques on Vietnamese Multi30k and UIT-ViIC.}
\label{tab:ablation}
\scriptsize
\setlength{\tabcolsep}{3pt} % giảm khoảng cách giữa các cột
\renewcommand{\arraystretch}{1.1}
\resizebox{\textwidth}{!}{
\begin{tabular}{|l|ccc|ccc|ccc|ccc|}
\hline
\multirow{2}{*}{Method} 
& \multicolumn{6}{c|}{Vietnamese Multi30k} 
& \multicolumn{6}{c|}{UIT-ViIC} \\
\cline{2-13}
& \multicolumn{3}{c|}{i2t} 
& \multicolumn{3}{c|}{t2i}
& \multicolumn{3}{c|}{i2t} 
& \multicolumn{3}{c|}{t2i} \\
\cline{2-13}
& R@1 & R@5 & R@10 
& R@1 & R@5 & R@10
& R@1 & R@5 & R@10 
& R@1 & R@5 & R@10 \\
\hline
Baseline 
& 49.3 & 80.1 & 86.7 
& 49.9 & 80.4 & 87.7 
& 20.2 & 44.8 & 59.9 
& 23.8 & 54.1 & 72.3 \\
\hline
+ Random Deletion
& 50.8 & 80.3 & 87.3
& 51.0 & 80.5 & 87.9
& 21.1 & 47.4 & 62.3
& 23.4 & 55.4 & 71.0 \\
\hline
+ Random Insertion
& 50.1 & 79.5 & 87.6
& 50.6 & 80.4 & 88.3
& 22.3 & 47.2 & 62.8
& 25.5 & 53.7 & 69.7 \\
\hline
+ Random Swap
& 49.6 & 79.4 & 87.5
& 49.9 & 80.8 & 88.2
& 21.7 & 46.6 & 62.2
& 23.4 & 52.4 & 72.0 \\
\hline
+ Back-Translation
& 50.1 & 79.5 & 87.7
& 50.9 & 80.1 & 87.8
& 21.1 & 46.7 & 63.0
& 24.2 & 54.4 & 72.3 \\
\hline
+ All Aug.
& \textbf{51.0} & \textbf{80.5} & \textbf{87.6}
& \textbf{51.2} & \textbf{80.7} & \textbf{88.9}
& \textbf{22.3} & \textbf{46.9} & \textbf{63.1}
& \textbf{23.4} & \textbf{55.0} & \textbf{72.7} \\
\hline
\end{tabular}
}
\end{table*}


\subsection{Effect of Text Augmentation Across Datasets}

Across both Vietnamese Multi30k and UIT-ViIC, augmentation consistently improves Recall@K and MeanR. The gains are particularly pronounced at stricter metrics such as Recall@1, indicating improved fine-grained alignment rather than merely increased similarity at coarse ranks. These findings confirm that data-centric strategies are effective for low-resource vision--language adaptation, especially when large-scale multilingual pretraining is unavailable.

\subsection{Qualitative Analysis}
To better understand the impact of data augmentation on retrieval performance, we present several qualitative examples comparing  the baseline model and the our method model. In the successful cases, the augmented model retrieves the correct image  or caption at rank 1 while the baseline model fails. These examples often  involve Vietnamese captions with complex modifier structures or passive constructions. We also present representative loss cases where augmentation degrades retrieval performance, suggesting that while augmentation improves robustness in many cases, it may introduce noise in visually ambiguous sscenes.


As shown in Table \ref{tab:win}, the augmented model successfully retrieves the correct results while the baseline model fails. For example, in the first case the Vietnamese caption contains a complex modifier structure describing the relationship between two people. The baseline model ranks the correct caption lower, while the augmented model retrieves the  correct result at rank 1. These examples suggest that augmentation improves the model's ability to handle complex Vietnamese descriptions.

\begin{table}[t]
\centering
\caption{Qualitative improvement examples for Image-to-Text (i2t) and Text-to-Image (t2i) retrieval.}
\label{tab:win}
\scriptsize
\setlength{\tabcolsep}{6pt}
\renewcommand{\arraystretch}{1.2}

\begin{tabular}{|c|>{\centering\arraybackslash}m{5cm}|>{\centering\arraybackslash}m{5.0cm}|>{\centering\arraybackslash}m{1.0cm}|>{\centering\arraybackslash}m{1.0cm}|}
\hline
\textbf{Task} & \textbf{Query} & \textbf{Ground Truth} & \textbf{Rank of baseline} & \textbf{Rank of our method} \\
\hline

i2t &
\includegraphics[width=1.0cm]{image/880.jpg} &
Một cô gái đeo mặt nạ đi trên vai một người đàn ông qua một lối đi đông đúc.(in EN: a girl wearing a mask rides on a man's shoulders through a crowded sidewalk) &
4 & 1 \\
\hline

t2i &
Một người phụ nữ đang ngồi trên một tảng đá rất lớn, mỉm cười nhìn vào máy ảnh với những cái cây ở phía sau.(in EN: a woman sitting on a very large rock smiling at the camera with trees in the background) &
\includegraphics[width=2.0cm]{image/868.jpg} &
8 & 1 \\
\hline
\end{tabular}
\end{table}

\begin{table}[t]
\centering
\caption{Qualitative failure examples for Image-to-Text (i2t) and Text-to-Image (t2i) retrieval.}
\label{tab:loss}
\scriptsize
\setlength{\tabcolsep}{5pt}
\renewcommand{\arraystretch}{1.15}
\begin{tabular}{|c|>{\centering\arraybackslash}m{5cm}|>{\centering\arraybackslash}m{5.0cm}|>{\centering\arraybackslash}m{1.0cm}|>{\centering\arraybackslash}m{1.0cm}|}
\hline
\textbf{Task} & \textbf{Query} & \textbf{Ground Truth} & \textbf{Rank of baseline} & \textbf{Rank of our method} \\
\hline

i2t &
\includegraphics[width=1.0cm]{image/306.jpg} &
Các cầu thủ bóng rổ bắn cho một bàn thắng trong một trận đấu.(in EN: basketball players shoot for a goal during a game) &
11 & 495 \\
\hline

t2i &
Một người đàn ông mặc áo da đen đang đứng trước một tấm biển \ 
(in EN: a man in a black leather coat is standing in front of a sign) &
\includegraphics[width=1.0cm]{image/669.jpg} &
25 & 145 \\
\hline
\end{tabular}
\end{table}


Table \ref{tab:loss} shows representative failure cases. Although augmentation improves many retrieval examples, it may also introduce noise in visually complex scenes. This observation indicates that augmentation does not universally improve performance and may occasionally degrade retrieval accuracy.

\subsection{Impact of Parameter-Efficient Fine-Tuning}
\label{sec:ft-impact}

% \begin{figure}[t]
%   \centering

%   % -------- (a) Image-to-Text --------
%   \begin{subfigure}{\linewidth}
%     \centering
%     \includegraphics[width=1.2\linewidth]{FJCAI2026-IEEE-conference-template/qualitative_i2t_rank1_5_10.png}
%     \caption*{(a) Image-to-Text}
%   \end{subfigure}

%   \vspace{4mm}

%   % -------- (b) Text-to-Image --------
%   \begin{subfigure}{\linewidth}
%     \centering
%     \includegraphics[width=1.0\linewidth]{FJCAI2026-IEEE-conference-template/qualitative_t2i_rank1_5_10.png}
%     \caption*{(b) Text-to-Image}
%   \end{subfigure}

%   \caption{Qualitative cross-modal retrieval results on the Vietnamese Multi30k dataset.}
%   \label{fig:qualitative}
% \end{figure}
\begin{figure}[t]
\centering

% ================= (a) =================
\textbf{\small (a) Image-to-Text}\par
\vspace{2mm}

\noindent
\begin{minipage}[t]{0.35\linewidth}
  \centering
  \vspace{0pt}
  \includegraphics[width=\linewidth]{i2t_query.png}
\end{minipage}\hfill
\begin{minipage}[t]{0.60\linewidth}
  \vspace{0pt}
  {\scriptsize
  \textbf{Rank 1: 0.482}\\
  VI: một cô gái mặc đồng phục karate phá vỡ một cái gậy với một cú đá phía trước.\\
  \textit{EN: a girl in karate uniform breaking a stick with a front kick.}\\[2mm]

  \textbf{Rank 5: 0.411}\\
  VI: Hai người và một con bò đeo dây ở một ngôi nhà.\\
  \textit{EN: two people and a leashed cow at a home.}\\[2mm]

  \textbf{Rank 10: 0.327}\\
  VI: Một người phụ nữ sử dụng một cái khoan trong khi một người đàn ông khác chụp ảnh cô ấy.\\
  \textit{EN: a woman uses a drill while another man takes her picture.}
  }
\end{minipage}

\vspace{4mm}

% ================= (b) =================
\textbf{\small (b) Text-to-Image}\par
\vspace{2mm}

{\scriptsize
\textbf{Query:} một cô gái mặc đồng phục karate phá vỡ một cái gậy với một cú đá phía trước.\\
\textit{EN: a girl in karate uniform breaking a stick with a front kick.}
}

\vspace{2mm}

{\scriptsize
\textbf{Rank 1 | 0.482}\hspace{10mm}
\textbf{Rank 5 | 0.348}\hspace{10mm}
\textbf{Rank 10 | 0.299}
}\par

\vspace{2mm}

\includegraphics[width=0.75\linewidth]{t2i_rank_1_5_10.png}

\vspace{2mm}
\caption{\footnotesize Qualitative cross-modal retrieval examples on the Vietnamese UIT-ViIC dataset.}
\label{fig:qualitative}
\end{figure}


We analyze the effect of parameter-efficient fine-tuning on the shared embedding space during Vietnamese adaptation. Since the CLIP image encoder remains frozen throughout training, the improvements mainly arise from adapting the multilingual text pathway and its lightweight projection layers. The visual representations learned during large-scale English-centric pretraining capture objects, actions, and spatial relationships that remain applicable across languages. As visual semantics are largely language-independent, cross-lingual discrepancies primarily stem from textual representations rather than visual features. The consistent gains in Recall@K and MeanR on both Vietnamese Multi30k and UIT-ViIC indicate that strengthening textual supervision alone can effectively enhance Vietnamese cross-modal alignment without modifying the pretrained visual backbone. Moreover, larger improvements on stricter metrics such as Recall@1 suggest improved fine-grained semantic discrimination rather than merely increased similarity.


To further evaluate the effectiveness of projection-layer adaptation, we compare it with a LoRA-based parameter-efficient fine-tuning strategy. In the LoRA setting, low-rank trainable matrices are injected into the attention projections of the multilingual text encoder while the backbone parameters remain frozen. Table~\ref{tab:lora_comparison} reports the retrieval performance on both Vietnamese Multi30k and UIT-ViIC. The LoRA-based model achieves slightly lower Recall@K scores, with a gap of around 2--3 points on Recall@1 compared with the projection-layer approach.

\begin{table*}[ht]
\centering
\caption{Comparison between projection-layer tuning and LoRA for cross-modal retrieval on Vietnamese Multi30k and UIT-ViIC.}
\label{tab:lora_comparison}
\scriptsize
\setlength{\tabcolsep}{3pt}
\renewcommand{\arraystretch}{1.1}
\resizebox{\textwidth}{!}{
\begin{tabular}{|l|ccc|ccc|ccc|ccc|}
\hline
\multirow{2}{*}{Method} 
& \multicolumn{6}{c|}{Vietnamese Multi30k} 
& \multicolumn{6}{c|}{UIT-ViIC} \\
\cline{2-13}

& \multicolumn{3}{c|}{i2t} 
& \multicolumn{3}{c|}{t2i}
& \multicolumn{3}{c|}{i2t} 
& \multicolumn{3}{c|}{t2i} \\
\cline{2-13}

& R@1 & R@5 & R@10 
& R@1 & R@5 & R@10
& R@1 & R@5 & R@10 
& R@1 & R@5 & R@10 \\
\hline

Projection (mCLIP)
& \textbf{49.3} & \textbf{80.1} & \textbf{86.7}
& \textbf{49.9} & \textbf{80.4} & \textbf{87.7}
& \textbf{20.2} & \textbf{44.8} & \textbf{59.9}
& \textbf{23.8} & \textbf{54.1} & \textbf{72.3} \\
\hline

LoRA
& 46.9 & 77.1 & 85.3
& 48.4 & 78.1 & 86.3
& 17.7 & 43.2 & 59.6
& 20.3 & 46.8 & 68.4 \\
\hline

\end{tabular}
}
\end{table*}


A comparison with prior benchmarks further highlights the effectiveness of our approach.
On the IGLUE benchmark, a benchmark introduced in the mCLIP study \cite{chen-etal-2023-mclip}, we evaluate the cross-modal retrieval capability of the proposed model. The benchmark has been used to assess multilingual multimodal models across approximately ten languages, including Vietnamese, providing a suitable setting to examine cross-lingual image–text alignment beyond English., mCLIP reports a Vietnamese Recall@1 of 19.4 for cross-modal retrieval \cite{radford2021learningtransferablevisualmodels}. Although evaluation settings differ, this score provides a useful reference.

Despite UIT-ViIC being substantially smaller, fine-tuning on native Vietnamese captions improves image-to-text Recall@1 to 20.2, exceeding the IGLUE result.
With text augmentation, Recall@1 further increases to 22.3, demonstrating that linguistic diversification yields consistent gains even under low-resource conditions.



Figure \ref{fig:qualitative} illustrates qualitative retrieval examples which consistent with these quantitative findings. Augmentation improves robustness to paraphrasing, flexible word order, and colloquial expressions in both translated Multi30k captions and native UIT-ViIC captions. Overall, these results confirm that text data augmentation plays a central role in effective Vietnamese adaptation, improving robustness and generalization without requiring large-scale manually curated image--text corpora. Since accurate retrieval alignment forms the foundation of dual-encoder vision--language systems, strengthening the shared embedding space is a prerequisite for extending the model to higher-level multimodal tasks such as visual question answering and grounded reasoning. The consistent improvements observed here therefore establish a solid basis for future expansion beyond retrieval.


\section{CONCLUSIONS}
\label{sec:conclusion}

This paper presents a data-efficient approach for adapting multilingual vision--language models to Vietnamese by combining parameter-efficient fine-tuning with targeted text data augmentation. By preserving pretrained encoders and enriching target-language supervision through controlled augmentation strategies, the proposed method strengthens cross-modal alignment under low-resource and domain-constrained conditions.

Experimental results on both translated and native Vietnamese datasets demonstrate consistent improvements in image--text retrieval performance. These findings confirm that data-centric strategies provide a practical and reproducible pathway for multilingual adaptation within parameter-efficient constraints. The gains observed across both synthetic and human-written supervision settings further indicate that carefully designed augmentation can enhance robustness to linguistic variation without requiring architectural modification. Future work will extend the evaluation of the adapted model to broader and more diverse Vietnamese domains to assess generalizability, multimodal benchmarks, including Vietnamese visual question answering and academic reasoning tasks, in order to assess its generalization beyond retrieval-oriented alignment.  Additional out-of-distribution evaluation on a more general Vietnamese dataset further confirms that the adapted model maintains reasonable retrieval performance beyond the specific-domain bias of UIT-ViIC.



%=========================================================================

\section*{ACKNOWLEDGMENT}


The authors sincerely thank the anonymous reviewers for their insightful comments and constructive suggestions, which significantly contributed to improving the quality and clarity of this manuscript. This research is funded by the Ministry of Science and Technology of Viet Nam (MOST) under grant number KC.01.15/21-30, entitled “Research and building of a scalable and unified platform for distributed video processing to serve smart urban management”.


%=========================================================================

%=========================================================================
% References
% BibTeX users please use
\bibliographystyle{IEEEtran} % sorted IEEE style
%\bibliographystyle{spmpsci_unsort}
\bibliography{template.bib} % name your BibTeX database
% \hfill {\it Received on October 14, 2003}

% \hfill {\it Accepted on January 14, 2004}
\end{document}
